\documentclass[11pt]{article}
\usepackage{amsmath,amsthm,amssymb}
\usepackage[margin=1in]{geometry}
\usepackage{enumitem}
\usepackage{hyperref}

\newtheorem{theorem}{Theorem}
\newtheorem{lemma}[theorem]{Lemma}
\newtheorem{proposition}[theorem]{Proposition}
\newtheorem{corollary}[theorem]{Corollary}
\theoremstyle{definition}
\newtheorem{definition}[theorem]{Definition}
\newtheorem{remark}[theorem]{Remark}

\title{The same-row part dies:\\
       structural proof for $\lambda = (3, 3, 1^q)$, $q \ge 2$}
\author{Clio}
\date{13 May 2026 (night)}

\begin{document}
\maketitle

\begin{abstract}
The r=1 inductive-reduction paper of 13 May
(\texttt{2026-05-13-r1-inductive-reduction.tex}) isolated a clean
conjecture (the ``same-row-dies conjecture'') whose verification
upgrades the r=1 reduction formula to an unconditional statement
on every test case. The conjecture was verified computationally
across $\lambda = (3, 3, 1^q)$ for $q \in \{2, 3, 4, 5\}$ but lacked a
structural proof. We supply that proof here for the family
$\lambda = (3, 3, 1^q)$, $q \ge 2$.

\textbf{Theorem.} For $\lambda = (3, 3, 1^q)$ with $q \ge 2$,
$n = q+6$, $\ell = q+2$, and $S$ the same-row part of
$E_{n-1}^+|_{V_\lambda}$ (the span of seminormal vectors $v_T$ for
SYTs of $\lambda$ with letters $n-1, n$ at the same-row pair
$(2, 2), (2, 3)$),
$$R'_{\ell+1}\,R'_{\ell+2}\,\cdots\,R'_{n-1}\,S \;=\; 0.$$

\textbf{Corollary.} The conditional reduction formula (Theorem 9 of
the r=1 paper) for $\lambda = (3, 3, 1^q)$ at $q \ge 2$ is now
unconditional:
\[
   B_{\ell+1}^{(\lambda)} V_\lambda
   \;=\;
   (T_\ell{+}1)(T_{\ell+1}{+}1)(T_{n-2}{+}1)\,
   \Phi_*\!\left(B^{(\mu_1)}_{\ell(\mu_1)+1} V_{\mu_1}\right).
\]
Combined with the May-12 evening result on $\mu_1 = (3, 2, 1^{q-1})$
(rigorous at $q \ge 4$, computational at $q = 2, 3$), this gives
$\dim B^{(\lambda)} = 2$ structurally at $q \ge 4$.

The proof uses three ingredients: (i) the
$H_q(S_{n-2})$-iso $S \cong V_{(3, 1^{q+1})}$ from
Lemma~3.3 of the r=1 paper; (ii) a clean operator commutation
identity inside $H_q(S_n)$ via index-gap-$\ge 2$ commutations;
(iii) the May-12 hook j-formula applied to the sub-shape
$\nu = (3, 1^{q+1})$ at sharp index. The non-obvious step is (ii):
combining $R'_{n-2}$ and $R'_{n-1}$ rearranges into
$(T_{n-3}+1)(T_{n-2}+1) R'_{n-3} R'_{n-2}$, exposing two
applications of $R'$-on-$S$ that match the sharp iteration on $\nu$.
\end{abstract}

\section{Setup and notation}\label{sec:setup}

Notation follows the r=1 paper (\texttt{2026-05-13-r1-inductive-reduction.tex}):
$H_q(S_n)$ is the Iwahori--Hecke algebra over $\mathbb{Q}(q)$;
$V_\lambda$ is the seminormal Specht module;
$E_j^\pm := E_j^\pm(V_\lambda)$ are the $\pm$-eigenspaces of $T_j$
(eigenvalues $q, -1$);
\[
   R'_k \;:=\; (T_{k-1}{+}1)(T_{k-2}{+}1)\cdots(T_1{+}1),
   \qquad
   B_j^{(\lambda)} V_\lambda \;:=\; R'_j R'_{j+1}\cdots R'_n V_\lambda.
\]

Throughout this paper $\lambda = (3, 3, 1^q)$ with $q \ge 2$,
$n = q + 6$, $\ell = q + 2$. The corners of $\lambda$ are
$c_1 = (2, 3)$ (length-preserving) and $c_* = (q+2, 1)$ (column-1
bottom); $\mu_1 := \lambda \setminus \{c_1, c_*\} = (3, 2, 1^{q-1})$.

The {\bf same-row pair} of $\lambda$ is the pair of cells
$(2, 2), (2, 3)$ at row 2 (this is the unique same-row pair of
$\lambda$, because $\lambda_2 = 3$ and $\lambda_3 = 1 \le 1 = \lambda_2 - 2$).

\begin{definition}[Same-row part $S$]\label{def:S}
$S \subseteq V_\lambda$ is the span of $\{v_T : T \in \mathrm{SYT}(\lambda),
\ T(2, 2) = n-1,\ T(2, 3) = n\}$. The remaining letters
$\{1, \ldots, n-2\}$ fill the cells of the sub-partition
\[
   \nu \;:=\; \lambda \setminus \{(2, 2), (2, 3)\}
        \;=\; (3, 1^{q+1}),
\]
on $n - 2 = q + 4$ letters. In particular $\dim S = f^\nu$.
\end{definition}

\begin{remark}
The Lemma~3.3 of the r=1 paper gives an explicit
$H_q(S_{n-2})$-equivariant iso $\Psi : S \xrightarrow{\sim} V_\nu$
(under the subalgebra $\langle T_1, \ldots, T_{n-3}\rangle$): $\Psi$
sends $v_T$ to $v_{T|_\nu}$, where $T|_\nu$ is the restriction of $T$
to letters $\{1, \ldots, n-2\}$. Equivariance: for $j \le n-3$, both
sides act by Hoefsmit's rule on positions of $j, j+1$, which lie in
$\nu$ (since $j+1 \le n-2$); hence the actions match.
\end{remark}

\section{The key commutation identity}\label{sec:commute}

\begin{lemma}[Commutation identity]\label{lem:commute}
Inside $H_q(S_n)$,
\begin{equation}\label{eq:commute}
   R'_{n-2}\, R'_{n-1}
   \;=\; (T_{n-3}+1)\,(T_{n-2}+1)\, R'_{n-3}\, R'_{n-2}.
\end{equation}
\end{lemma}

\begin{proof}
First, $R'_{n-1} = (T_{n-2}+1)\, R'_{n-2}$ by definition. Hence
\[
   R'_{n-2}\, R'_{n-1} \;=\; R'_{n-2}\,(T_{n-2}+1)\, R'_{n-2}.
\]
Next, $R'_{n-2} = (T_{n-3}+1)\, R'_{n-3}$, and $R'_{n-3}$ is a product of
factors $(T_j+1)$ with $j \le n - 4$. Each such $T_j$ has $|j - (n-2)| \ge 2$,
so $T_j$ commutes with $T_{n-2}$; equivalently $R'_{n-3}$ commutes with
$(T_{n-2}+1)$. Therefore
\[
   R'_{n-2}\,(T_{n-2}+1)
   \;=\; (T_{n-3}+1)\, R'_{n-3}\,(T_{n-2}+1)
   \;=\; (T_{n-3}+1)\,(T_{n-2}+1)\, R'_{n-3}.
\]
Substituting:
\[
   R'_{n-2}\, R'_{n-1}
   \;=\; (T_{n-3}+1)\,(T_{n-2}+1)\, R'_{n-3}\, R'_{n-2}. \qedhere
\]
\end{proof}

\begin{remark}
The identity~\eqref{eq:commute} is shape-independent — it holds inside
$H_q(S_n)$ as an operator equation, independent of the chosen
representation.
\end{remark}

\section{The hook j-formula on $\nu = (3, 1^{q+1})$}\label{sec:hook}

\begin{lemma}[May-12 hook j-formula]\label{lem:hook-j}
For $\nu = (3, 1^{q+1})$ on $m = q+4$ letters with $q \ge 2$
(so $m \ge 6$),
\begin{equation}\label{eq:hook-j}
   B^{(\nu)}_{m-1} V_\nu \;:=\; R'^{(\nu)}_{m-1}\, R'^{(\nu)}_m\, V_\nu
   \;\subseteq\; E_1^-|_{V_\nu}.
\end{equation}
\end{lemma}

\begin{proof}
This is the j-formula for the hook $(3, 1^{m-3})$, proved
structurally in \texttt{2026-05-12-hook-j-formula.tex} (commit
\texttt{9b532b6}). The proof tracks the image of $R'_{n-1} R'_n$
step by step in the seminormal basis: Phase A ($R'_m$) gives
dim $m - 2 = q + 2$ with a basis of $T_{m-1}$-2-block
$+q$-eigenvectors plus the same-row $+q$-eigenvector at row 1;
Phase B ($R'_{m-1}$) starts with $(T_1+1)$ collapsing to dim~1,
then a shifting-window argument shows the surviving support is
a single SYT $v_{T^*}$ with $T^*(1, 1) = 1$ on letters in
$\{n-2k+3, \ldots, n\}$ (top range), giving the support in
$E_1^-$ via min element $\ge 3 \ge 2$. The full Phase B argument
lands $B^{(\nu)}_{m-1} V_\nu$ inside $E_1^-$.
\end{proof}

\begin{remark}
The full rank-zero theorem $\Pi^{S_m}|_{V_\nu} = 0$ at every $m \ge 6$
follows from~\eqref{eq:hook-j} by one further $R'_{m-2}$
application: $R'_{m-2}$ leads (on the right, applied first) with
$(T_1+1)$, which kills $E_1^-$. We use only~\eqref{eq:hook-j}
here, not the full rank-zero theorem.
\end{remark}

\section{Main theorem}\label{sec:main}

\begin{theorem}[Same-row dies for $(3, 3, 1^q)$]\label{thm:main}
For $\lambda = (3, 3, 1^q)$ with $q \ge 2$, $n = q+6$,
$\ell = q+2$, and $S$ the same-row part of
Definition~\ref{def:S},
\begin{equation}\label{eq:main}
   R'_{\ell+1}\,R'_{\ell+2}\,\cdots\,R'_{n-1}\,S \;=\; 0.
\end{equation}
\end{theorem}

\begin{proof}
Since $\ell + 1 = q + 3 = n - 3$, the iteration in~\eqref{eq:main}
is exactly three $R'$-applications:
\[
   R'_{\ell+1}\, R'_{\ell+2}\, R'_{n-1}
   \;=\; R'_{n-3}\, R'_{n-2}\, R'_{n-1}.
\]
We show this acts as zero on $S$, in three steps.

\paragraph{Step 1.} \emph{Apply the commutation
identity~\eqref{eq:commute}.} We have
\begin{equation}\label{eq:step1}
   R'_{n-3}\, R'_{n-2}\, R'_{n-1}
   \;=\; R'_{n-3}\, \big((T_{n-3}+1)\,(T_{n-2}+1)\,R'_{n-3}\,R'_{n-2}\big)
   \;=\; R'_{n-3}\,(T_{n-3}+1)\,(T_{n-2}+1)\,R'_{n-3}\,R'_{n-2}.
\end{equation}

\paragraph{Step 2.} \emph{$R'_{n-3}\, R'_{n-2}|_S$ lies in
$E_1^-|_{V_\lambda}$.}

Restrict to $S$ and use the iso $\Psi: S \xrightarrow{\sim} V_\nu$
(Lemma~3.3 of the r=1 paper), which is $T_j$-equivariant for
$j \le n-3$. The operators $R'_{n-3}, R'_{n-2}$ are products of
factors $(T_j+1)$ with $j \le n-3$, so each is in the
$\langle T_1, \ldots, T_{n-3}\rangle$-subalgebra and respects
$\Psi$:
\[
   \Psi(R'_{n-3}|_S \cdot R'_{n-2}|_S \cdot v)
   \;=\; R'^{(\nu)}_{n-3}\, R'^{(\nu)}_{n-2}\, \Psi(v)
   \quad \text{for all } v \in S.
\]
Now $R'^{(\nu)}_{n-2} = R'^{(\nu)}_{m}$ and $R'^{(\nu)}_{n-3} = R'^{(\nu)}_{m-1}$
(since $m = n - 2$). So
\[
   \Psi(R'_{n-3}|_S \cdot R'_{n-2}|_S \cdot V_\nu)
   \;=\; R'^{(\nu)}_{m-1}\, R'^{(\nu)}_m\, V_\nu
   \;=\; B^{(\nu)}_{m-1} V_\nu.
\]
By the May-12 hook j-formula (Lemma~\ref{lem:hook-j}),
$B^{(\nu)}_{m-1} V_\nu \subseteq E_1^-|_{V_\nu}$.

Pulling back via $\Psi^{-1}$ (which is $T_1$-equivariant since $1 \le n-3$):
\[
   R'_{n-3}\, R'_{n-2}\, S
   \;\subseteq\; \Psi^{-1}(E_1^-|_{V_\nu})
   \;=\; E_1^-|_S
   \;\subseteq\; E_1^-|_{V_\lambda}.
\]
(The last inclusion uses that $T_1$ acts on $S \subseteq V_\lambda$
by the same operator whether viewed in $S$ or in $V_\lambda$,
so $E_1^-|_S = S \cap E_1^-|_{V_\lambda}$.)

\paragraph{Step 3.} \emph{$(T_{n-3}+1)\,(T_{n-2}+1)$ preserves
$E_1^-|_{V_\lambda}$.}

For $n \ge 5$, both $T_{n-3}$ and $T_{n-2}$ have index gap $\ge 2$
from $T_1$, so they commute with $T_1$. Hence $(T_{n-3}+1)$ and
$(T_{n-2}+1)$ commute with $T_1$ and preserve the $T_1$-eigenspaces,
in particular $E_1^-|_{V_\lambda}$.

Combining with Step 2,
\[
   (T_{n-3}+1)\,(T_{n-2}+1)\, R'_{n-3}\, R'_{n-2}\, S
   \;\subseteq\; (T_{n-3}+1)\,(T_{n-2}+1)\, E_1^-|_{V_\lambda}
   \;\subseteq\; E_1^-|_{V_\lambda}.
\]

\paragraph{Step 4.} \emph{$R'_{n-3}$ kills $E_1^-|_{V_\lambda}$.}

$R'_{n-3} = (T_{n-4}+1)(T_{n-5}+1)\cdots(T_2+1)(T_1+1)$. The
rightmost factor is $(T_1+1)$, which is applied first to any vector.
On $E_1^-$, the operator $T_1$ acts as $-1$, so $(T_1+1)|_{E_1^-} = 0$.
Hence $R'_{n-3}\, E_1^-|_{V_\lambda} = 0$.

\paragraph{Conclusion.} By~\eqref{eq:step1}, Step 3, and Step 4,
\[
   R'_{n-3}\, R'_{n-2}\, R'_{n-1}\, S
   \;=\; R'_{n-3}\big( (T_{n-3}+1)(T_{n-2}+1)\, R'_{n-3}\, R'_{n-2}\, S \big)
   \;\subseteq\; R'_{n-3}\, E_1^-|_{V_\lambda} \;=\; 0. \qedhere
\]
\end{proof}

\section{Consequences}\label{sec:cons}

\begin{corollary}[Reduction formula unconditional for $(3, 3, 1^q)$, $q \ge 2$]\label{cor:reduction}
For $\lambda = (3, 3, 1^q)$ with $q \ge 2$,
\begin{equation}\label{eq:reduction}
   B_{\ell+1}^{(\lambda)} V_\lambda
   \;=\;
   (T_\ell{+}1)(T_{\ell+1}{+}1)(T_{n-2}{+}1)\,
   \Phi_*\!\left(B_{\ell(\mu_1)+1}^{(\mu_1)} V_{\mu_1}\right).
\end{equation}
\end{corollary}

\begin{proof}
By Lemma~4.1 of the r=1 paper (the $\Phi_*$-piece commutation,
unconditional) and the Phase A decomposition
$E_{n-1}^+|_{V_\lambda} = \Phi_*(V_{\mu_1}) \oplus S$ (Lemma~3.4 of the
r=1 paper), it suffices to verify $R'_{\ell+1}\cdots R'_{n-1}\,S = 0$.
This is Theorem~\ref{thm:main}.
\end{proof}

\begin{corollary}[Structural upper bound $\dim B \le 2$ for $(3, 3, 1^q)$, $q \ge 4$]\label{cor:dim}
For $\lambda = (3, 3, 1^q)$ with $q \ge 4$,
$\dim B_{\ell+1}^{(\lambda)} V_\lambda \le 2$, structurally.
\end{corollary}

\begin{proof}
By Corollary~\ref{cor:reduction},
$\dim B^{(\lambda)} \le \dim B^{(\mu_1)}$ (since the outer factors and
$\Phi_*$ are injective on the inner space). For $q \ge 4$,
$|\mu_1| = q + 4 \ge 8$, so the May-12 evening theorem applies:
$\dim B^{(\mu_1)} = 2$. Hence $\dim B^{(\lambda)} \le 2$.
\end{proof}

\begin{remark}[Lower bound]
The matching lower bound $\dim B^{(\lambda)} \ge 2$ is verified
computationally for $q \in \{2, 3, 4, 5\}$ (Theorem~12 of the r=1
paper). A structural lower bound via chain decomposition + outer
factor injectivity remains open --- the proof sketch in
\texttt{2026-05-13-dim2-331/proof-sketch.md} attempts this but
identifies a position-of-$(n-2)$ invariant that fails to be
preserved by the outer factor $(T_{n-2}+1)$. The structural
lower bound is not addressed in this paper.
\end{remark}

\section{Computational verification}\label{sec:comp}

Mod-$P$ verification at $P = 100\,003$, $q = 11$ (the seed
specialisation), for $\lambda = (3, 3, 1^q)$ at $q \in \{2, 3, 4\}$:

\begin{itemize}[leftmargin=2em,topsep=0.3em,itemsep=0.2em]
\item Step A (operator identity~\eqref{eq:commute}): \texttt{True} at
$n \in \{8, 9, 10\}$.
\item Step 2 ($R'_{n-3}\,R'_{n-2}\,S \subseteq E_1^-$): rank of
$(T_1+1) \cdot R'_{n-3}\,R'_{n-2}\,S = 0$ at $n \in \{8, 9, 10\}$
(verified directly).
\item Step 3 ($(T_{n-3}+1)(T_{n-2}+1)$ preserves $E_1^-$): rank of
$(T_1+1)$ applied to $(T_{n-3}+1)(T_{n-2}+1)\,R'_{n-3}\,R'_{n-2}\,S$
is also $0$ at $n \in \{8, 9, 10\}$.
\item Final image $R'_{\ell+1}\,R'_{\ell+2}\,R'_{n-1}\,S = 0$:
direct rank computation gives dim~$0$ at $n \in \{8, 9, 10, 11\}$.
\end{itemize}

Scripts: \texttt{\textasciitilde/projects/scratch/2026-05-13-same-row-dies/}.

The decay pattern (Phase A) on $S$ matches the rank-zero iteration
on $\nu = (3, 1^{q+1})$:
\[
   \dim S = f^\nu
   \;\xrightarrow{R'_{n-1}}\; q + 2
   \;\xrightarrow{R'_{n-2}}\; 1
   \;\xrightarrow{R'_{n-3}}\; 0,
\]
with the same dimension drops as $V_\nu \xrightarrow{R'_m} q+2
\xrightarrow{R'_{m-1}} 1 \xrightarrow{R'_{m-2}} 0$ for $\nu$ on
$m = q + 4$ letters --- exactly the parallel
predicted by the iso $S \cong V_\nu$ + Phase A on $\nu$.

\section{Where the proof generalises}\label{sec:general}

The proof of Theorem~\ref{thm:main} uses three ingredients, all
shape-sensitive:

\begin{enumerate}[leftmargin=2em,topsep=0.3em,itemsep=0.2em]
\item \textbf{The $H_q(S_{n-2})$-iso} $S \cong V_\nu$, where
$\nu$ is the sub-shape obtained by removing the same-row pair. This
is a general fact: it holds whenever $\lambda$ has a same-row pair
(Lemma~3.3 of the r=1 paper).
\item \textbf{The commutation identity~\eqref{eq:commute}}. This is
purely operator-theoretic in $H_q(S_n)$ and holds independent of the
representation, so it transfers to any $\lambda$.
\item \textbf{The j-formula for $\nu$ at sharp index}:
$B^{(\nu)}_{j_0(\nu)} V_\nu \subseteq E_1^-|_{V_\nu}$. This is the
shape-specific input. For $\nu = (3, 1^{q+1})$ a hook, it is May-12.
\end{enumerate}

Therefore the proof extends to {\bf any} r=1 shape $\lambda$ with a
same-row pair, provided:

\begin{enumerate}[label=(\roman*),topsep=0.2em,itemsep=0.1em]
\item the sub-shape $\nu$ has $\ell(\nu) = \ell(\lambda)$ (so that the
sharp index $j_0(\nu)$ matches the next $R'$-application in $\lambda$),
\item the j-formula $B^{(\nu)}_{j_0(\nu)} V_\nu \subseteq E_1^-|_{V_\nu}$
is known.
\end{enumerate}

\paragraph{Concrete other cases the proof handles directly.}
\begin{itemize}[leftmargin=2em,topsep=0.2em,itemsep=0.1em]
\item Hooks $(k, 1^{n-k})$ with $k \ge 3$: same-row at row 1,
$\nu = (k-2, 1^{n-k})$ is a hook with one shorter top row.
Condition~(i) holds. The j-formula on $\nu$: for $k = 3$,
$\nu = (1, 1^{n-3}) = (1^{n-2})$ is the sign rep, trivially
in $E_1^-$; for $k \ge 4$, $\nu = (k-2, 1^{n-k})$ is again a hook,
and one would need to combine with the May-12 hook proof or extend
it inductively. (However: hooks $(k, 1^{n-k})$ already have rank-zero
proved unconditionally for $k \le 3$ in May-12 and structurally for
all $k$ in the morning's chain decomposition --- so this corollary
of our proof is not new for hooks.)
\item $(k, k, 1^q)$ at $k = 3$: this is the family of this paper.
\item $(k, k, 1^q)$ at $k \ge 4$: $\nu = (k, k-2, 1^q)$ is not a
hook. The j-formula on $\nu$ is not yet structurally established for
all such shapes. Stays open under the current proof.
\end{itemize}

The cleanest next step in this thread is to prove the j-formula for
the family $\nu = (k, k-2, 1^q)$ (and more generally non-hook
non-r=2-corner shapes), which would propagate the same-row-dies
proof through to $(k, k, 1^*)$ at $k \ge 4$.

\section{Honest status}\label{sec:honest}

\paragraph{What we proved.} Same-row-dies for
$\lambda = (3, 3, 1^q)$ at every $q \ge 2$, unconditionally, via the
$H_q(S_{n-2})$-iso $S \cong V_{(3, 1^{q+1})}$, an operator commutation,
and the May-12 hook j-formula.

\paragraph{What we did not prove.} The same-row-dies conjecture for
general r=1 shapes with same-row pairs (e.g., $(k, k, 1^*)$ at
$k \ge 4$). The obstruction is the j-formula on the non-hook
sub-shape $\nu$, which requires its own structural work.

\paragraph{What this closes.} The r=1 paper's Theorem~12 (dim $B = 2$
for $(3, 3, 1^q)$, $q \in \{2, 3, 4, 5\}$, conditional on same-row-dies)
is now unconditional at $q \ge 2$. The dim is structural at $q \ge 4$
via May-12 evening on $\mu_1 = (3, 2, 1^{q-1})$; structural plus
computational at $q = 2, 3$ pending a structural dim-2 lower bound
for $\mu_1 = (3, 2)$ and $(3, 2, 1)$.

\end{document}
